\section{Перечень процессов}\label{PROCESSES.List}

Управление сертификатами конечных участников ИОК реализуется через следующие
процессы.

\begin{enumerate}
\item
\code{Enroll}~--- 
выпуск сертификата для стороны, которая располагает действительным
удостоверением, но не обязательно действительным сертификатом.
\item
\code{Reenroll}~--- 
обновление действительного сертификата.
\item
\code{Spawn}~--- 
выпуск нового сертификата для стороны, которая располагает действительным
сертификатом.
\item
\code{Retrieve}~--- 
получение сертификата, который был запрошен в процессах~\code{Enroll},
\code{Reenroll}, \code{Spawn} и выпуск которого задерживается.
\item
\code{Setpwd}~--- 
установка (изменение) пароля отзыва сертификата.
\item
\code{Revoke}~--- 
отзыв сертификата.
\end{enumerate}

Каждый процесс представляет собой последовательность определенных процедур. 
При ошибке в любой из процедур процесс также завершается с ошибкой.

В каждом процессе обязательно участвует УЦ~--- ИУЦ или ПУЦ. Процесс
включает процедуры отправки запроса УЦ и получения соответствующего ответа.
Сторона, отправляющая запрос, должна располагать сертификатом УЦ.

Форматы запросов и ответов схематически представлены в
таблице~\ref{Table.PROCESSES.Fmt}. Используются имена типов данных и их
компонентов, описанные в разделе~\ref{FMT}. Нижний индекс~C означает субъекта
сертификата, нижний индекс~опРЦ~--- оператора~РЦ, нижний индекс~агУЦ~---
агента~УЦ.
%
Нижний индекс у контейнера~\code{SignedData} указывает на подписанта, у
контейнера~\code{EnvelopedData}~--- на получателя конвертованных данных, у
других объектов~--- на их владельцев.

\begin{table}[bht]
\caption{Схемы форматов запросов/ответов}
\label{Table.PROCESSES.Fmt}
\begin{tabular}{|l|l|}
\hline
\multicolumn{2}{|c|}{Процесс \code|Enroll|}\\
\hline
\hline
\rule{0pt}{18pt}
Запрос &
$\code|EnvelopedData|_{\text{УЦ}}\bigl(
\code|SignedData|_{\text{РЦ~| опРЦ}}
(\code|CertificateRequest|_{\text{С}})\bigr)$ или\\
&
$\code|EnvelopedData|_{\text{УЦ}}(
\code|CertificateRequest|_{\text{С}}[\code|enrollPwd|])$\\[6pt]
\hline                                      
%
\rule{0pt}{18pt}
Ответ &
$\code|EnvelopedData|_{\text{РЦ~| опРЦ~| С}}(\code|Certificate|_{\text{С}})$ или\\
&
$\code|SignedData|_{\text{агУЦ}}(\code|BPKIResp|)$\\[6pt]
\hline                                     
\hline
\multicolumn{2}{|c|}{Процессы \code|Reenroll| и \code|Spawn|}\\
\hline
\hline
\rule{0pt}{15pt}
Запрос &
$\code|EnvelopedData|_{\text{УЦ}}\bigl(
\code|SignedData|_{\text{С}}
(\code|CertificateRequest|_{\text{С}})\bigr)$\\[3pt]
\hline                                      
%
\rule{0pt}{15pt}
Ответ &
$\code|EnvelopedData|_{\text{С}}(\code|Certificate|_{\text{С}})$ или\\
&
$\code|SignedData|_{\text{агУЦ}}(\code|BPKIResp|)$\\[3pt]
\hline                                     
\hline
\multicolumn{2}{|c|}{Процесс \code|Retrieve|}\\
\hline
\hline
\rule{0pt}{15pt}
Запрос &
\code|BPKIRetrieveReq|\\[3pt]
\hline                                      
%
\rule{0pt}{15pt}
Ответ &
$\code|EnvelopedData|_{\text{РЦ~| опРЦ~| С}}(\code|Certificate|_{\text{С}})$ или\\
&
$\code|SignedData|_{\text{агУЦ}}(\code|BPKIResp|)$\\[3pt]
\hline                                     
\hline
\multicolumn{2}{|c|}{Процесс \code|Setpwd|}\\
\hline
\hline
\rule{0pt}{15pt}
Запрос &
$\code|EnvelopedData|_{\text{УЦ}}\bigl(
\code|SignedData|_{\text{C}}(\code|revokePwd|)\bigr)$\\[3pt]
\hline                                      
%
\rule{0pt}{15pt}
Ответ &
$\code|SignedData|_{\text{агУЦ}}(\code|BPKIResp|)$\\[3pt]
\hline                                     
\hline
\multicolumn{2}{|c|}{Процесс \code|Revoke|}\\
\hline
\hline
\rule{0pt}{15pt}
Запрос &
$\code|EnvelopedData|_{\text{УЦ}}
\bigl(\code|SignedData|_{\text{С}}(\code|BPKIRevoke|)\bigr)$ или\\
&
$\code|EnvelopedData|_{\text{УЦ}}(\code|BPKIRevoke|)$\\[3pt]
\hline                                      
%
\rule{0pt}{15pt}
Ответ &
$\code|SignedData|_{\text{агУЦ}}(\code|BPKIResp|)$\\[3pt]
\hline                                     
\end{tabular}
\end{table}

Перечисленные в таблице запросы и ответы транспортируются в виде пакетов 
HTTP. Правила транспорта определяются в~\ref{TRANSPORT}.

